\documentclass[11pt]{article}
\usepackage{graphicx}
\setlength{\parindent}{0pt}
\usepackage{hyperref}
\usepackage{enumitem}
\usepackage[utf8]{inputenc} 
\usepackage[T1,T2A]{fontenc}
\usepackage[brazil]{babel}
\usepackage{lipsum}
\usepackage[left=1.06cm,top=1.7cm,right=1.06cm,bottom=0.49cm]{geometry}
\usepackage{fontawesome}

\usepackage{hyphenat}
\hyphenpenalty=10000 % полностью запрещает переносы

\begin{document}
\begin{center}
    \textbf{Denis Boborukhin}\\ 
    \hrulefill
\end{center}

\newcommand{\infofield}[2]{\mbox{{\normalsize #1}\hspace{0.5em}#2}}
\newcommand{\email}[1]{\infofield{\faEnvelope}{\href{mailto:#1}{#1}}}
\newcommand{\location}[1]{\infofield{\faMapMarker}{#1}}
\newcommand{\linkedin}[2]{\infofield{\faLinkedin}{\href{#2}{#1}}}
\newcommand{\github}[2]{\infofield{\faGithub}{\href{#2}{#1}}}
\newcommand{\telegram}[2]{\infofield{\faSend}{\href{#2}{#1}}}

\begin{center}
    \location{Moscow, Russia}
    \email{denis.boborukhin@gmail.com}
    \telegram{@dboborukhin}{https://t.me/dboborukhin} \\
    \linkedin{Denis Boborukhin}{https://www.linkedin.com/in/denis-boborukhin/}
    \github{@denisboborukhin}{https://github.com/denisboborukhin}
\end{center}

\vspace{0.5pt}

\begin{center}
    \textbf{Education}
\end{center}

\textbf{Skolkovo Institute of Science and Technology: Skoltech} \hfill Moscow

Master of Engineering, Engineering Systems: Product Lifecycle Management \hfill 2025-2027
\begin{itemize}[noitemsep, topsep=0pt, partopsep=0pt, parsep=0pt]
    \item Systems architecture, product lifecycle management 
    \item Technological entrepreneurship
\end{itemize}

\vspace{12pt}

\textbf{Moscow Institute of Physics and Technology: MIPT} \hfill Moscow

Bachelor, Radio Engineering and Cybernetics \hfill 2021-2025
\begin{itemize}[noitemsep, topsep=0pt, partopsep=0pt, parsep=0pt]
    \item The use and application of the C/C++ languages, system programming 
    \item Architecture of computing cores and memory subsystems of modern microprocessors
\end{itemize}

\vspace{12pt}

\begin{center}
    \textbf{Experience}
\end{center}
\textbf{Yadro} \hfill Moscow

\textbf{Software Engineer} \hfill July 2023 – Present
\begin{itemize}[noitemsep, topsep=0pt, partopsep=0pt, parsep=0pt]
    \item Developed QEMU-based profiling tools for metric collection, enabling application analysis to identify performance
     bottlenecks and generate insights for potential compiler and microarchitectural optimizations.
    \item Researched and developed methods for detecting microarchitectural performance bottlenecks through \\
     instruction-level analysis, providing recommendations for design improvements.
    \item Developed a tool for selecting small, representative application sections that accurately capture the overall workload,
     enabling performance analysis on high-fidelity simulators.
    \item Implemented a system for metric storage and analytics, providing insights into applications behavior
     and enabling visualization of large-scale performance data.
\end{itemize}

\begin{center}
    \textbf{Activities}
\end{center}

\textbf{FidgetSense (Partner)} \hfill September 2025 – Present

Develops a hardware-software system of three wireless smart rings that continuously track finger movements
 to detect anxiety-related fidgeting. Algorithms analyze movement dynamics
 to quantify anxiety levels and evaluate treatment effectiveness through observed dynamic over time.

Achievements:
\begin{itemize}[noitemsep, topsep=0pt, partopsep=0pt, parsep=0pt]
    \item Triple Point Social Tech Competition: Third place
    \item Skoltech Innovation Workshop (Medical Cohort): Third place
\end{itemize}

\begin{center}
    \textbf{Skills \& Interests}
\end{center}

\textbf{Technical:} C/C++, Python, Linux, ML, architecture of computing cores and memory subsystem, Go, MBSE

\textbf{Language:} Russian(native), English(B2)

\textbf{Interests:} state of the art technologies, systems engineering, system programming, system modeling \& simulation


\end{document}
